%% Chapter 8: The Existential Theory of the Reals
%% Status: skeleton -- learning goals and section outline fixed, prose to be written.

\chapter{The Existential Theory of the Reals}
\label{ch:etr}

A striking number of natural problems --- in geometry, in graph drawing, in game
theory, in machine learning --- are not known to be in \textbf{NP}, and the
obstruction is always the same: a solution is a tuple of real numbers that may
need exponentially many bits to write down. The class $\ER$ is the right home
for these problems.

\begin{goals}
  \poolgoal{8.1}{give a definition of the algorithmic problem \ETR{} and of the complexity class $\ER$ based on it}
  \poolgoal{8.2}{give the machine model description of $\ER$ and a sketch of a proof that it is equivalent with the logic-based definition}
  \poolgoal{8.3}{show that $\mathrm{NP} \subseteq \ER$}
  \poolgoal{8.4}{define the order types and the order type problem, and give a motivation to study order types}
  \poolgoal{8.5}{define \textsc{ETR-AM} and show that it is $\ER$-complete}
  \poolgoal{8.6}{define the partial order type problem and show that it is $\ER$-complete}
\end{goals}

\section{\texorpdfstring{The problem \ETR{}}{The problem ETR}}
\label{sec:etr-problem}

\todo{Write this section.}

\section{\texorpdfstring{The class $\ER$}{The class exists-R}}
\label{sec:er-class}

\todo{Write this section.}
% Planned content: Logic-based definition, the machine model over the reals, and the sandwich $\mathrm{NP} \subseteq \ER \subseteq \mathrm{PSPACE}$.

\section{\texorpdfstring{$\mathrm{NP} \subseteq \ER$}{NP is contained in exists-R}}
\label{sec:np-in-er}

\todo{Write this section.}

\section{Order types}
\label{sec:order-types}

\todo{Write this section.}
% Planned content: Definition, the order type realizability problem, and why order types are the right combinatorial abstraction of a point set.

\section{Normal forms: \textsc{ETR-INV} and \textsc{ETR-AM}}
\label{sec:etr-am}

\todo{Write this section.}
% Planned content: Restricting the variable range and the shape of the equations, without losing $\ER$-hardness.

\section{The partial order type problem}
\label{sec:partial-order-types}

\todo{Write this section.}

\section{Notes and further reading}
\label{sec:etr-further-reading}

This chapter is based on lecture notes by Tillmann Miltzow.

\todo{Add the literature: Shor and Canny for $\ER \subseteq \mathrm{PSPACE}$,
Mn\"ev's universality theorem, Richter-Gebert's work on order types, and
Schaefer's surveys of $\ER$-completeness.}

\section{Exercises}
\label{sec:ex-etr}

\todo{Write this section.}
